\documentclass[11pt,a4paper]{article}
\usepackage[margin=1.9cm,top=1.7cm]{geometry}
\usepackage{graphicx,array,booktabs,tabularx,amssymb,xcolor,enumitem}
\usepackage[hidelinks]{hyperref}
\usepackage{fancyhdr}
\pagestyle{fancy}\fancyhf{}
\lhead{\small 34870 Loudspeaker project -- listening test, operator sheet}\rhead{\small Group 10, System D}\cfoot{\small\thepage}
\setlength{\parskip}{3pt}\setlength{\parindent}{0pt}
\newlist{chk}{itemize}{1}
\setlist[chk]{label=$\square$,itemsep=1pt,leftmargin=1.6em}
\newcommand{\warn}[1]{\textcolor{red!70!black}{\textbf{#1}}}
\begin{document}

{\LARGE\bfseries Listening test: operator sheet}\\[2pt]
{\large Readiness checklist, procedure and sample orders}\\[2pt]
{\small Keep this sheet \textbf{outside} the listening room: it holds the sample orders.\\
Listener sheet: \texttt{ListeningTest\_Questionnaire.pdf}.\\
Background: \texttt{lab-day/results\_2026-11-12/LabDay\_12Nov\_Results.pdf}.}

\section*{What the test is for}
The anechoic results say the speaker is flat to 5.5\,dB (150\,Hz--15\,kHz) once the MiniDSP is set right.
The listening test answers three questions the chamber cannot:
\begin{enumerate}[itemsep=0pt]
\item \textbf{Part A:} how do listeners rate the final design, and does anything stand out (bass, the 1.5--3\,kHz plateau, treble)?
\item \textbf{Part B:} is the woofer level right \emph{in this room with this amplifier}? Blind comparison of Out 1 at $-7.8$, $-5.8$ and $-3.8$\,dB.
\item \textbf{Part C:} is the tweeter level right? Blind comparison of Rp $=8.2\,\Omega$ (now) and $3.3\,\Omega$ (before 12 Nov).
\end{enumerate}

\section{Readiness checklist}
\subsection*{Bring}
\begin{chk}
\item The passive board with \textbf{Rp $=8.2\,\Omega$} fitted, plus a \textbf{3.3\,$\Omega$} resistor ($\geq5$\,W, course basket) for Part C, and a screwdriver.
\item The vent plug (one vent open at 160\,mm, the other plugged).
\item USB stick \texttt{34870\_G10\_ListeningTest\_2026-11-18/}: MiniDSP filter files, check tones, the four test excerpts, the vinyl albums, these PDFs.
\item Mads's PC with the \textbf{Schiit Modi} DAC, its USB cable, and an RCA cable (or RCA to the room's input). Play the FLAC files from the stick, no streaming, no Bluetooth; system volume at 100\,\% and player volume fixed (bit-perfect if the player allows).
\item Multimeter (AC volts) for the electrical check. Tape and a pen to mark the volume knob.
\item Printed: one questionnaire per listener plus 3 spares, this sheet, pens, a stopwatch (phone). Three chairs, with seat labels Left / Middle / Right.
\item A thin cloth to hang in front of the speaker and the board, if one is available (not required).
\end{chk}

\subsection*{Music: four excerpts, 90\,s each, the same files for every listener}
\begin{tabularx}{\linewidth}{@{}l>{\raggedright}p{5.0cm}>{\raggedright\arraybackslash}X@{}}\toprule
Track & What it tests & Excerpt (from Mads's vinyl rips, folder \texttt{4\_Test\_excerpts})\\\midrule
1 female voice & midrange naturalness, the 1.5--3\,kHz plateau & Fugees -- \emph{Ready or Not}\\
2 male voice & colouration, ``boxiness'', sibilance & Michael Bubl\'e -- \emph{Have Yourself a Merry Little Christmas}\\
3 bass-rich & bass amount and boom (crossover at 325\,Hz) & Miles Davis -- \emph{So What}\\
4 full band & treble, cymbals, clarity & Electric Light Orchestra -- \emph{Turn to Stone}\\\bottomrule
\end{tabularx}

The excerpts are cut, faded and levelled to the same loudness ($-23$\,LUFS) by \texttt{make\_excerpts.py}, and they are
\textbf{dual mono} (left = right), so MiniDSP routing stays Input 1 $\rightarrow$ Outputs 1 and 2 only. Listen to all four once
before the test and fix any start time in \texttt{make\_excerpts.py} if an excerpt starts in the wrong place.

\subsection*{Signal chain}
\textbf{PC $\rightarrow$ USB $\rightarrow$ Schiit Modi $\rightarrow$ RCA $\rightarrow$ MiniDSP input $\rightarrow$ MiniDSP outputs 1/2 $\rightarrow$ amplifier
ch 1/2 $\rightarrow$ room pass-through $\rightarrow$ woofer / board.} \warn{Never connect the DAC straight to the amplifier:} without the
MiniDSP there is no crossover, and the midrange (which has only a low-pass on the board) gets the full bass at full power.
The Modi puts out about 2\,V; start with the PC volume low and watch the MiniDSP input level for clipping.

\subsection*{Set up in the IEC room}
\begin{chk}
\item Box at the room's marked loudspeaker position, tweeter at ear height of a seated listener if possible; \textbf{three chairs}: Middle at the marked listening position (on axis), Left and Right beside it, as close as is comfortable, all facing the box. Use the same three chair positions in every session.
\item Board inside the room, wired as in the chamber: amplifier channel 1 $\rightarrow$ woofer, channel 2 $\rightarrow$ board $\rightarrow$ mid + tweeter. Red to + on every driver.
\item MiniDSP (the room's own card): routing only Input 1 (and 2) $\rightarrow$ Output 1 and Output 2. PEQ, Comp, FIR \textbf{bypassed}; only Xover active.
\item Output 1 Xover $\rightarrow$ Advanced: delete everything, paste \texttt{MiniDSP\_woofer\_out1.txt}, Process. Output 2: the same with \texttt{MiniDSP\_midtw\_out2.txt}.
\item \warn{Output 2 must hold exactly two biquads and no other high-pass.} On 12 Nov it held the filter twice (a $-11$\,dB hole at 400\,Hz). Check the basic Xover page too.
\item Gains: \textbf{Out 1 $=-5.8$\,dB, Out 2 $=0$\,dB}. Invert off, delay 0, sample rate 96\,kHz. Save the configuration.
\item Amplifier tone controls flat (bass and treble centred), loudness OFF, balance centred.
\end{chk}

\subsection*{Electrical check (5 minutes, volume low)}
Play each check tone (\texttt{tone\_0100Hz}, \texttt{tone\_0325Hz}, \texttt{tone\_1000Hz}, $-12$\,dBFS) and read the AC
voltage at the two amplifier outputs with the speaker connected. Only ratios matter, so any volume setting works.

\begin{tabular}{@{}l>{\centering}p{3.5cm}>{\centering}p{3.5cm}>{\raggedright\arraybackslash}p{5.2cm}@{}}\toprule
Tone & Ch 1 (woofer) & Ch 2 (mid + tw) & Expected\\\midrule
100\,Hz & $V_1$ = \rule{1.5cm}{0.4pt} & \rule{1.5cm}{0.4pt} & ch 2 close to 0\\
325\,Hz & \rule{1.5cm}{0.4pt} & \rule{1.5cm}{0.4pt} & both $=0.50\times$ their own reference. \warn{Ch 2 at $0.25\times$ = filter doubled}\\
1\,kHz & \rule{1.5cm}{0.4pt} & $V_2$ = \rule{1.5cm}{0.4pt} & ch 1 close to 0\\\bottomrule
\end{tabular}

Also write down $V_1/V_2$ = \rule{1.5cm}{0.4pt}. It is 0.51 ($-5.8$\,dB) if the room amplifier's two channels have equal gain.
\textbf{By ear:} at 100\,Hz only the woofer plays, at 1\,kHz only the midrange (and the tweeter faintly). If it is the other way round,
the amplifier channels are swapped (as in runs S and S\_2 on 12 Nov).

\subsection*{Playback level}
\begin{chk}
\item With Out 1 at $-5.8$\,dB, play track 4 and set a comfortably loud level at the middle chair (about 75--80\,dB(A) if a sound level meter or phone app is at hand).
\item Mark the volume knob with tape. \warn{Never touch the volume again during the test.} Write the setting here: \rule{3cm}{0.4pt}
\end{chk}

\subsection*{Phone form}
\begin{minipage}{0.2\linewidth}\includegraphics[width=\linewidth]{qr_form.png}\end{minipage}\hfill
\begin{minipage}{0.76\linewidth}
Listeners can answer on their phone instead of paper: \textbf{listening.madsrudolph.dev} (same questions, one screen at a time,
saved even if the page reloads). Each listener enters the session number and their seat. Results and CSV: the private results link
(key in \texttt{web/.admin-key}); it fills in the blind conditions automatically from the table in Section~4.
\end{minipage}

\section{Roles}
\begin{tabularx}{\linewidth}{@{}lX@{}}\toprule
Operator (outside) & plays the excerpts, changes Out 1 between samples, swaps Rp in Part C, holds the order table. Never talks to the listeners about the settings.\\
Guide (inside) & reads the script, hands out the sheet, keeps time. \textbf{Does not know the order}, so cannot give anything away (double blind).\\
Listeners & \textbf{three per session} (Left, Middle, Right), two sessions or more; people outside the group if possible. Group members may listen too: they tick it on their sheet, analyse them separately.\\\bottomrule
\end{tabularx}
Agree on a signal between guide and operator before the first session (a text message ``next'' works; keep the door closed).

\section{Procedure per session: three listeners at once (about 35 minutes)}
All three hear the same playback, so the sample order belongs to the \textbf{session}, not to the person.
\begin{enumerate}[itemsep=2pt]
\item \textbf{Before they enter:} Out 1 at \textbf{$-5.8$\,dB}, Rp \textbf{8.2\,$\Omega$} (Part A plays the design).
The guide tells the three listeners the \textbf{session number}; each writes it and their seat (Left / Middle / Right) on the sheet or the phone form.
The operator looks up that session's order in the table below.
\item \textbf{Guide reads aloud:} ``You will hear one loudspeaker playing music excerpts. Please stay in your seat, keep your head still and face
the speaker. Don't talk, don't react, and don't look at each other's answers. Nobody in this room knows the settings. There are no right
answers; rate what you hear. The scales are explained at the top of the sheet.'' Let them read page 1.
\item \textbf{Part A} (design settings): tracks 1, 2, 3, 4, each once, with about 40\,s between tracks to fill in the column.
\item \textbf{Part B}: four samples. Each sample = track 1 then track 3. Between samples the operator sets Out 1 to the next
value in the session's order (Part B table) and waits 40\,s, until all three have rated. Same order of tracks every time.
\item Put Out 1 back to $-5.8$\,dB.
\item \textbf{Part C}: the guide takes \textbf{all three listeners} outside, away from the door. The operator goes in and sets Rp for sample X
(Part C table), even if it is already right: always unscrew and refit, so the sound of a swap gives nothing away. Listeners and guide
return; track 4 then track 2; rate X. Out again, the operator sets the Y value, back in, track 4 then track 2, rate Y.
Leave Rp at \textbf{8.2\,$\Omega$} afterwards.
\item Last question, thank them, collect the sheets (check that each has the session number and a seat). The operator fills in the grey line at the bottom of page 2.
\end{enumerate}
Each session takes about 35 minutes (Part A 9, Part B 15, Part C 9, intro 2) for three listeners: two sessions (6 listeners) need about 1 h 15 min. Short on time: stop each Part B and C sample after the first track, or run Part C with only some listeners; Parts A and B matter most.

\section{Sample orders}
Part B conditions (MiniDSP Output 1 gain): \textbf{L} $=-7.8$\,dB (bass $-2$\,dB), \textbf{D} $=-5.8$\,dB (design),
\textbf{H} $=-3.8$\,dB (bass $+2$\,dB). D is played twice as a hidden repeat: if a listener's two D ratings differ by more than 3 points,
that listener is guessing. Part C: \textbf{R8} $=$ Rp 8.2\,$\Omega$, \textbf{R3} $=$ Rp 3.3\,$\Omega$.

\begin{center}
\renewcommand{\arraystretch}{1.25}
\begin{tabular}{@{}c|cccc|cc|l@{}}\toprule
Session & \multicolumn{4}{c|}{Part B, sample 1 -- 4 (Out 1 gain)} & \multicolumn{2}{c|}{Part C} & Time / notes\\
no. & 1 & 2 & 3 & 4 & X & Y & \\\midrule
1 & D $-5.8$ & L $-7.8$ & D $-5.8$ & H $-3.8$ & R8 & R3 & \rule{4cm}{0.4pt}\\
2 & H $-3.8$ & D $-5.8$ & L $-7.8$ & D $-5.8$ & R3 & R8 & \rule{4cm}{0.4pt}\\
3 & D $-5.8$ & H $-3.8$ & L $-7.8$ & D $-5.8$ & R8 & R3 & \rule{4cm}{0.4pt}\\
4 & L $-7.8$ & D $-5.8$ & H $-3.8$ & D $-5.8$ & R3 & R8 & \rule{4cm}{0.4pt}\\
5 & D $-5.8$ & H $-3.8$ & D $-5.8$ & L $-7.8$ & R8 & R3 & \rule{4cm}{0.4pt}\\
6 & D $-5.8$ & L $-7.8$ & H $-3.8$ & D $-5.8$ & R3 & R8 & \rule{4cm}{0.4pt}\\\midrule
7 & \multicolumn{6}{l|}{as session 1 \quad(8 = 2, 9 = 3, \dots)} & \rule{4cm}{0.4pt}\\\bottomrule
\end{tabular}
\end{center}
Number the sessions 1, 2, 3, \dots\ in order, and do not show this table to anyone inside. Each listener is identified by session + seat (e.g.\ 2M); the results page uses the same.

\section{Afterwards}
\begin{chk}
\item Type every sheet into \texttt{listening-test/results.csv} (columns in \texttt{results\_template.csv}): one row per answer, with the condition taken from the table above.
\item For the report: mean and spread per attribute in Part A; in Parts B and C, mean overall score and mean ``amount'' per condition.
\item Decide: if ``bass amount'' for D averages near 0 and D scores at least as well as L and H, keep Out 1 $=-5.8$\,dB. If H wins and its bass amount is near 0, move to $-3.8$\,dB, and so on. The same rule for Rp in Part C.
\item Whatever is changed: measure it in the chamber (one \texttt{S} run), so the report shows the final setting as a curve.
\item Leave the room as found: MiniDSP configuration saved, board and plug back with the box, Rp at 8.2\,$\Omega$.
\end{chk}

\end{document}
